\chapter{User manual}
\label{app:manual}

\section{Requirements}
Linux or macOS-like shell; \code{g++} with C++17, \code{cmake} $\geq$ 3.16, \code{python3} $\geq$ 3.11 with \code{venv}, \code{node} $\geq$ 18 with \code{npm}, and \code{gcc} (used as assembler and linker for the native back end). For the SQL differential test no extra install is needed: Python's standard \code{sqlite3} is the oracle.

\section{Getting started}
\begin{lstlisting}[language=bash,style=tiny]
git clone <repository> nova && cd nova
./scripts/dev.sh              # builds the compiler, installs dependencies, starts backend + studio
# open the printed URL (default http://localhost:5173)
./scripts/serve.sh            # alternative: one process serving the built studio on http://localhost:8000
\end{lstlisting}

\section{Using the studio}
\begin{enumerate}
  \item Open \textbf{Compiler Playground}, pick an example, press \textbf{Compile} (Ctrl+Enter) or \textbf{Run} (Ctrl+Shift+Enter).
  \item Click the stages in the pipeline bar, or open \emph{Lexical}, \emph{Syntax}, \emph{AST}, \emph{Semantic}, \emph{Intermediate Repr.}, \emph{CFG \& Data-flow}, \emph{Optimization Lab}, \emph{Target Code (VM)} and \emph{Native x86-64} from the menu.
  \item In \emph{Optimization Lab} switch passes off and on and watch the rewrite history and the measured effect.
  \item In \emph{Debugger} set breakpoints by clicking line numbers; step forward and \emph{back}.
  \item Break a program on purpose and open \textbf{AI Assistant}: apply a verified fix, or press \emph{Auto-fix all}; for a runtime trap use \emph{Runtime analysis}.
  \item Open \textbf{SQL Compiler}, pick an example script, press \textbf{Run} (Ctrl+Enter) and inspect the \emph{plan}, \emph{optimizer}, \emph{bytecode} and \emph{ast} tabs.
\end{enumerate}

\section{Command-line use}
\begin{lstlisting}[language=bash,style=tiny]
novac=compiler/build/novac
$novac compile --run < examples/02_factorial.nova | jq .execution     # full JSON report
$novac compile --passes=constant_folding,dce < prog.nova               # choose passes
$novac asm < prog.nova                                                 # x86-64 assembly
scripts/nova-native.sh prog.nova -o prog && ./prog                     # native executable
$novac sql < examples/sql/01_orders_report.sql | jq '.statements[-1].rows'
$novac fuzz --count=1000 --seed=1                                      # differential fuzz
$novac grammar | jq .states                                            # LR(1) tables as JSON
\end{lstlisting}

\section{HTTP API examples}
\begin{lstlisting}[language=bash,style=tiny]
curl -s localhost:8000/api/compile -H 'content-type: application/json' \
     -d '{"source":"fn main(){ print(6*7); }","run":true}' | jq .execution.optimized.output
curl -s localhost:8000/api/sql -H 'content-type: application/json' \
     -d '{"script":"CREATE TABLE t(a INT); INSERT INTO t VALUES (1),(2); SELECT SUM(a) FROM t;"}'
curl -s localhost:8000/api/suggest -H 'content-type: application/json' \
     -d '{"source":"fn main(){ let x: int = 1.5; }"}' | jq '.fixes[0].title'
\end{lstlisting}

\section{Language server in an editor}
Run \code{scripts/nova-lsp.sh} as an LSP server on stdio for files with the extension \code{.nova}. A minimal VS Code client is in \code{editors/vscode} (run \code{npm install}, open the folder, press F5, and set \code{nova.serverScript} to the absolute path of the script).

\section{Optional: enabling the LLM assistant}
Export \code{ANTHROPIC\_API\_KEY} before starting the backend; \code{NOVA\_LLM\_MODEL} selects the model. Without a key the \emph{Ask Claude} tab reports that the model is not configured and everything else works.

\section{Running the verification}
See Listing~\ref{lst:verifycmds}. \code{scripts/test.sh} runs everything; \code{scripts/sanitize.sh} adds the sanitizer builds; \code{scripts/bench.sh} reproduces the performance table.

\section{Troubleshooting}
\begin{table}[H]\centering\small
\begin{tabularx}{\textwidth}{@{}p{0.34\textwidth}X@{}}\toprule
\textbf{Symptom} & \textbf{Fix}\\ \midrule
``Backend offline'' in the studio & start it with \code{scripts/dev.sh}; if ports are busy set \code{NOVA\_API\_PORT} and \code{NOVA\_WEB\_PORT}\\
\code{compiler binary not found} (HTTP 503) & run \code{scripts/build.sh}\\
Native build fails & install \code{gcc}; check that the target is x86-64 Linux\\
Native/Compiler Explorer C viewer unavailable & the viewer needs \code{gcc} or \code{clang} on \code{PATH}\\
\code{ModuleNotFoundError} in backend tests & use the project's virtual environment: \code{backend/.venv/bin/python -m pytest}\\ \bottomrule
\end{tabularx}
\end{table}

\chapter{Diagnostic catalogue}
Every diagnostic has a stable code, a phase, an exact source range, an explanation and, where possible, a suggestion. The table lists all \NumDiagCodes{} codes of the compilers (generated from the source). Codes starting with \code{W} are warnings; \code{E-ICE} marks an internal compiler error (a bug report, never expected).

\begingroup\small
\begin{longtable}{@{}llp{0.55\textwidth}@{}}\toprule
\textbf{Code} & \textbf{Phase} & \textbf{Category}\\ \midrule\endhead
\texttt{E-ICE-001} &  & Internal compiler error\\
\texttt{E-LEX-001} & lexical & Invalid character\\
\texttt{E-LEX-002} & lexical & Invalid operator\\
\texttt{E-LEX-003} & lexical & Unterminated comment\\
\texttt{E-LEX-004} & lexical & Malformed number\\
\texttt{E-LEX-005} & lexical & Malformed number\\
\texttt{E-LEX-006} & lexical & Integer literal too large\\
\texttt{E-LEX-007} & lexical & Unterminated string\\
\texttt{E-LEX-008} & lexical & Invalid escape\\
\texttt{E-SEM-001} & semantic & Undefined variable\\
\texttt{E-SEM-002} & semantic & Duplicate declaration\\
\texttt{E-SEM-003} & semantic & Duplicate declaration\\
\texttt{E-SEM-004} & semantic & Type mismatch\\
\texttt{E-SEM-005} & semantic & Type mismatch\\
\texttt{E-SEM-006} & semantic & Undefined function\\
\texttt{E-SEM-007} & semantic & Wrong argument count\\
\texttt{E-SEM-008} & semantic & Invalid return\\
\texttt{E-SEM-009} & semantic & Missing return\\
\texttt{E-SEM-010} & semantic & Missing entry point\\
\texttt{E-SEM-011} & semantic & Invalid entry point\\
\texttt{E-SEM-012} & semantic & Invalid use of function\\
\texttt{E-SEM-013} & semantic & Void value used\\
\texttt{E-SEM-014} & semantic & Type mismatch\\
\texttt{E-SEM-015} & semantic & Type mismatch\\
\texttt{E-SEM-016} & semantic & Type mismatch\\
\texttt{E-SEM-017} & semantic & Type mismatch\\
\texttt{E-SQL-001} & NOVA-SQL & Unterminated string\\
\texttt{E-SQL-002} & NOVA-SQL & Malformed number\\
\texttt{E-SQL-003} & NOVA-SQL & Integer too large\\
\texttt{E-SQL-004} & NOVA-SQL & Invalid character\\
\texttt{E-SQL-005} & NOVA-SQL & Unterminated comment\\
\texttt{E-SQL-010} & NOVA-SQL & Syntax error\\
\texttt{E-SQL-100} & NOVA-SQL & Unknown table\\
\texttt{E-SQL-101} & NOVA-SQL & Unknown column\\
\texttt{E-SQL-102} & NOVA-SQL & Ambiguous column\\
\texttt{E-SQL-103} & NOVA-SQL & Type mismatch\\
\texttt{E-SQL-104} & NOVA-SQL & Unknown function\\
\texttt{E-SQL-105} & NOVA-SQL & Wrong argument count\\
\texttt{E-SQL-106} & NOVA-SQL & Invalid aggregate\\
\texttt{E-SQL-107} & NOVA-SQL & Invalid grouping\\
\texttt{E-SQL-108} & NOVA-SQL & Duplicate table\\
\texttt{E-SQL-109} & NOVA-SQL & Duplicate column\\
\texttt{E-SQL-110} & NOVA-SQL & Wrong value count\\
\texttt{E-SQL-111} & NOVA-SQL & Constraint violation\\
\texttt{E-SQL-112} & NOVA-SQL & Invalid ORDER BY\\
\texttt{E-SQL-113} & NOVA-SQL & Resource limit\\
\texttt{E-SYN-001} & syntax & Missing semicolon\\
\texttt{E-SYN-002} & syntax & Missing parenthesis\\
\texttt{E-SYN-003} & syntax & Missing brace\\
\texttt{E-SYN-004} & syntax & Unexpected end of file\\
\texttt{E-SYN-005} & syntax & Invalid expression\\
\texttt{E-SYN-006} & syntax & Syntax error\\
\texttt{W-SEM-001} & semantic & Unused variable\\
\texttt{W-SEM-002} & semantic & Unreachable code\\
\bottomrule
\end{longtable}

\endgroup
\captionof{table}{All diagnostic codes.}

\chapter{Example programs and their output}
The \NumExamples{} programs shipped in \code{examples/}, each with the real compiler result: the program output for valid programs, the diagnostics for the error examples, and the runtime trap for the runtime-error examples.

\subsection*{01\_hello\_arithmetic.nova --- Arithmetic \& branching}
\noindent The introductory NOVA program: variables, expressions, if/else.
\begin{lstlisting}[language=nova,]
fn main() {
    let age: int = 22;
    let marks: float = 91.5;
    let passed: bool = true;

    let result: int = 10 + 20 * 2;

    if (result > 30) {
        print(result);
    } else {
        print(0);
    }
    print(marks);
    print(passed);
    print(age);
}
\end{lstlisting}
\begin{tcolorbox}[colback=green!4,colframe=green!45!black,title={Program output (optimized VM)},fonttitle=\bfseries\small,boxsep=2pt,left=4pt,right=4pt,top=2pt,bottom=2pt]
\begin{lstlisting}[basicstyle=\ttfamily\scriptsize,frame=none,backgroundcolor=\color{white}]
50
91.5
true
22
\end{lstlisting}
\end{tcolorbox}
\subsection*{02\_factorial.nova --- Recursive factorial}
\noindent Function calls, parameters, return values and recursion (watch call frames in the debugger).
\begin{lstlisting}[language=nova,]
fn fact(n: int) -> int {
    if (n <= 1) {
        return 1;
    }
    return n * fact(n - 1);
}

fn main() {
    let i: int = 1;
    while (i <= 6) {
        print(fact(i));
        i = i + 1;
    }
}
\end{lstlisting}
\begin{tcolorbox}[colback=green!4,colframe=green!45!black,title={Program output (optimized VM)},fonttitle=\bfseries\small,boxsep=2pt,left=4pt,right=4pt,top=2pt,bottom=2pt]
\begin{lstlisting}[basicstyle=\ttfamily\scriptsize,frame=none,backgroundcolor=\color{white}]
1
2
6
24
120
720
\end{lstlisting}
\end{tcolorbox}
\subsection*{03\_fibonacci.nova --- Fibonacci (iterative + recursive)}
\noindent Two implementations of the same function; both print the same values.
\begin{lstlisting}[language=nova,]
fn fib_rec(n: int) -> int {
    if (n < 2) { return n; }
    return fib_rec(n - 1) + fib_rec(n - 2);
}

fn fib_iter(n: int) -> int {
    let a: int = 0;
    let b: int = 1;
    let i: int = 0;
    while (i < n) {
        let next: int = a + b;
        a = b;
        b = next;
        i = i + 1;
    }
    return a;
}

fn main() {
    let n: int = 0;
    while (n < 10) {
        print(fib_iter(n));
        n = n + 1;
    }
    print(fib_rec(12));
}
\end{lstlisting}
\begin{tcolorbox}[colback=green!4,colframe=green!45!black,title={Program output (optimized VM)},fonttitle=\bfseries\small,boxsep=2pt,left=4pt,right=4pt,top=2pt,bottom=2pt]
\begin{lstlisting}[basicstyle=\ttfamily\scriptsize,frame=none,backgroundcolor=\color{white}]
0
1
1
2
3
5
8
13
21
34
144
\end{lstlisting}
\end{tcolorbox}
\subsection*{04\_optimization\_demo.nova --- Optimization showcase}
\noindent Constant folding, propagation, CSE, algebraic identities, a dead branch and dead stores.
\begin{lstlisting}[language=nova,]
fn main() {
    let base: int = 4 * 8 + 2;          // folds to 34
    let scale: int = 2;
    let x: int = base * scale + 0;      // x*1/+0 identities
    let y: int = base * scale;          // common subexpression with x
    let unused: int = 7 * 6;            // dead store
    let debug: bool = false;

    if (debug) {                        // branch is never taken
        print(unused);
    }

    let t: int = (x + y) * 2;           // strength reduction
    print(t);
    print(x - x);
    print(y % 1);
}
\end{lstlisting}
\begin{tcolorbox}[colback=green!4,colframe=green!45!black,title={Program output (optimized VM)},fonttitle=\bfseries\small,boxsep=2pt,left=4pt,right=4pt,top=2pt,bottom=2pt]
\begin{lstlisting}[basicstyle=\ttfamily\scriptsize,frame=none,backgroundcolor=\color{white}]
272
0
0
\end{lstlisting}
\end{tcolorbox}
\subsection*{05\_primes.nova --- Prime sieve by trial division}
\noindent Nested while loops, modulo and boolean logic.
\begin{lstlisting}[language=nova,]
fn is_prime(n: int) -> bool {
    if (n < 2) { return false; }
    let d: int = 2;
    while (d * d <= n) {
        if (n % d == 0) { return false; }
        d = d + 1;
    }
    return true;
}

fn main() {
    let n: int = 2;
    let count: int = 0;
    while (n < 40) {
        if (is_prime(n)) {
            print(n);
            count = count + 1;
        }
        n = n + 1;
    }
    print(count);
}
\end{lstlisting}
\begin{tcolorbox}[colback=green!4,colframe=green!45!black,title={Program output (optimized VM)},fonttitle=\bfseries\small,boxsep=2pt,left=4pt,right=4pt,top=2pt,bottom=2pt]
\begin{lstlisting}[basicstyle=\ttfamily\scriptsize,frame=none,backgroundcolor=\color{white}]
2
3
5
7
11
13
17
19
23
29
31
37
...
\end{lstlisting}
\end{tcolorbox}
\subsection*{06\_floats\_and\_casts.nova --- Floats and explicit casts}
\noindent NOVA has no implicit conversions: use to\_float / to\_int.
\begin{lstlisting}[language=nova,]
fn average(a: int, b: int, c: int) -> float {
    return to_float(a + b + c) / 3.0;
}

fn main() {
    let avg: float = average(90, 85, 77);
    print(avg);
    let rounded: int = to_int(avg * 100.0);
    print(rounded);
    print(avg > 84.0 && avg < 85.0);
}
\end{lstlisting}
\begin{tcolorbox}[colback=green!4,colframe=green!45!black,title={Program output (optimized VM)},fonttitle=\bfseries\small,boxsep=2pt,left=4pt,right=4pt,top=2pt,bottom=2pt]
\begin{lstlisting}[basicstyle=\ttfamily\scriptsize,frame=none,backgroundcolor=\color{white}]
84.0
8400
false
\end{lstlisting}
\end{tcolorbox}
\subsection*{07\_scopes.nova --- Nested scopes \& shadowing}
\noindent Inner blocks may shadow outer names; see the scope tree in Semantic Analysis.
\begin{lstlisting}[language=nova,]
fn main() {
    let x: int = 1;
    {
        let x: int = 2;
        print(x);
        {
            let x: int = 3;
            print(x);
        }
        print(x);
    }
    print(x);
}
\end{lstlisting}
\begin{tcolorbox}[colback=green!4,colframe=green!45!black,title={Program output (optimized VM)},fonttitle=\bfseries\small,boxsep=2pt,left=4pt,right=4pt,top=2pt,bottom=2pt]
\begin{lstlisting}[basicstyle=\ttfamily\scriptsize,frame=none,backgroundcolor=\color{white}]
2
3
2
1
\end{lstlisting}
\end{tcolorbox}
\subsection*{08\_gcd.nova --- Euclid's GCD}
\noindent A classic loop with modulo; good for stepping in the debugger.
\begin{lstlisting}[language=nova,]
fn gcd(a: int, b: int) -> int {
    while (b != 0) {
        let r: int = a % b;
        a = b;
        b = r;
    }
    return a;
}

fn main() {
    print(gcd(48, 18));
    print(gcd(1071, 462));
    print(gcd(17, 5));
}
\end{lstlisting}
\begin{tcolorbox}[colback=green!4,colframe=green!45!black,title={Program output (optimized VM)},fonttitle=\bfseries\small,boxsep=2pt,left=4pt,right=4pt,top=2pt,bottom=2pt]
\begin{lstlisting}[basicstyle=\ttfamily\scriptsize,frame=none,backgroundcolor=\color{white}]
6
21
1
\end{lstlisting}
\end{tcolorbox}
\subsection*{09\_arrays.nova --- Arrays: bubble sort}
\noindent Heap-allocated int[] with literals, new, indexing, assignment and len(); arrays are passed by reference.
\begin{lstlisting}[language=nova,]
fn swap(a: int[], i: int, j: int) {
    let t: int = a[i];
    a[i] = a[j];
    a[j] = t;
}

fn sort(a: int[]) {
    let n: int = len(a);
    let i: int = 0;
    while (i < n) {
        let j: int = 0;
        while (j < n - 1 - i) {
            if (a[j] > a[j + 1]) { swap(a, j, j + 1); }
            j = j + 1;
        }
        i = i + 1;
    }
}

fn main() {
    let a: int[] = [5, 2, 9, 1, 7, 3];
    sort(a);
    let i: int = 0;
    while (i < len(a)) { print(a[i]); i = i + 1; }
    let squares: int[] = new int[5];
    i = 0;
    while (i < 5) { squares[i] = i * i; i = i + 1; }
    print(squares[4]);
}
\end{lstlisting}
\begin{tcolorbox}[colback=green!4,colframe=green!45!black,title={Program output (optimized VM)},fonttitle=\bfseries\small,boxsep=2pt,left=4pt,right=4pt,top=2pt,bottom=2pt]
\begin{lstlisting}[basicstyle=\ttfamily\scriptsize,frame=none,backgroundcolor=\color{white}]
1
2
3
5
7
9
16
\end{lstlisting}
\end{tcolorbox}
\subsection*{10\_strings.nova --- Strings}
\noindent Immutable string literals with concatenation, equality, len() and print.
\begin{lstlisting}[language=nova,]
fn greet(name: string) -> string {
    return "Hello, " + name + "!";
}

fn main() {
    let s: string = greet("NOVA");
    print(s);
    print(len(s));
    print(s == "Hello, NOVA!");
    print(s != greet("world"));
    let line: string = "";
    let i: int = 0;
    while (i < 3) { line = line + "ab"; i = i + 1; }
    print(line);
}
\end{lstlisting}
\begin{tcolorbox}[colback=green!4,colframe=green!45!black,title={Program output (optimized VM)},fonttitle=\bfseries\small,boxsep=2pt,left=4pt,right=4pt,top=2pt,bottom=2pt]
\begin{lstlisting}[basicstyle=\ttfamily\scriptsize,frame=none,backgroundcolor=\color{white}]
Hello, NOVA!
12
true
true
ababab
\end{lstlisting}
\end{tcolorbox}
\subsection*{e01\_undefined\_variable.nova --- Error: undefined variable}
\noindent Misspelled identifier; the compiler suggests a fix using edit distance.
\begin{lstlisting}[language=nova,]
fn main() {
    let counter: int = 10;
    print(countr + 1);
}
\end{lstlisting}
\begin{tcolorbox}[colback=red!4,colframe=red!60!black,title={Compiler output},fonttitle=\bfseries\small,boxsep=2pt,left=4pt,right=4pt,top=2pt,bottom=2pt]
\begin{lstlisting}[basicstyle=\ttfamily\scriptsize,frame=none,backgroundcolor=\color{white}]
E-SEM-001 at 6:11: Variable 'countr' is not declared
    hint: Did you mean 'counter'?
\end{lstlisting}
\end{tcolorbox}
\subsection*{e02\_missing\_semicolon.nova --- Error: missing semicolon}
\noindent Syntax error found by the LR(1) parser; location points just after the previous token.
\begin{lstlisting}[language=nova,]
fn main() {
    let a: int = 5
    print(a);
}
\end{lstlisting}
\begin{tcolorbox}[colback=red!4,colframe=red!60!black,title={Compiler output},fonttitle=\bfseries\small,boxsep=2pt,left=4pt,right=4pt,top=2pt,bottom=2pt]
\begin{lstlisting}[basicstyle=\ttfamily\scriptsize,frame=none,backgroundcolor=\color{white}]
E-SYN-001 at 5:19: Expected ';' but found 'print'
    hint: Insert ';' at the end of the previous statement.
\end{lstlisting}
\end{tcolorbox}
\subsection*{e03\_type\_mismatch.nova --- Error: type mismatch}
\noindent Mixing int and float, a float assigned to an int variable, and a non-bool condition.
\begin{lstlisting}[language=nova,]
fn main() {
    let a: int = 5;
    let b: float = 2.5;
    let c: float = a + b;
    let d: int = 3.5;
    if (a) {
        print(c);
    }
}
\end{lstlisting}
\begin{tcolorbox}[colback=red!4,colframe=red!60!black,title={Compiler output},fonttitle=\bfseries\small,boxsep=2pt,left=4pt,right=4pt,top=2pt,bottom=2pt]
\begin{lstlisting}[basicstyle=\ttfamily\scriptsize,frame=none,backgroundcolor=\color{white}]
E-SEM-005 at 7:20: Operator '+' cannot be applied to int and float
    hint: NOVA has no implicit int/float conversion. Wrap the int operand in to_float(...) or the float operand in to_int(...).
E-SEM-004 at 8:18: Cannot initialise 'd' of type int with a value of type float
    hint: Use to_int(<expr>) to truncate explicitly.
E-SEM-004 at 9:9: If condition must be bool but has type int
    hint: Compare explicitly, e.g. (x != 0).
\end{lstlisting}
\end{tcolorbox}
\subsection*{e04\_wrong\_arguments.nova --- Error: wrong function calls}
\noindent Wrong argument count, wrong argument type and a call to an undefined function.
\begin{lstlisting}[language=nova,]
fn add(a: int, b: int) -> int {
    return a + b;
}

fn main() {
    print(add(1));
    print(add(1, 2.0));
    print(addd(1, 2));
}
\end{lstlisting}
\begin{tcolorbox}[colback=red!4,colframe=red!60!black,title={Compiler output},fonttitle=\bfseries\small,boxsep=2pt,left=4pt,right=4pt,top=2pt,bottom=2pt]
\begin{lstlisting}[basicstyle=\ttfamily\scriptsize,frame=none,backgroundcolor=\color{white}]
E-SEM-007 at 9:11: Function 'add' expects 2 argument(s) but 1 were given
    hint: Add the missing argument(s).
E-SEM-004 at 10:18: Argument 2 of 'add' has type float but parameter 'b' is int
    hint: Pass a int value.
E-SEM-006 at 11:11: Function 'addd' is not defined
    hint: Did you mean 'add'?
\end{lstlisting}
\end{tcolorbox}
\subsection*{e05\_duplicate\_and\_return.nova --- Error: duplicate declaration \& missing return}
\noindent Redeclaring a variable in the same scope, and a function that can fall off its end.
\begin{lstlisting}[language=nova,]
fn sign(n: int) -> int {
    if (n > 0) {
        return 1;
    }
}

fn main() {
    let x: int = 1;
    let x: int = 2;
    print(sign(x));
}
\end{lstlisting}
\begin{tcolorbox}[colback=red!4,colframe=red!60!black,title={Compiler output},fonttitle=\bfseries\small,boxsep=2pt,left=4pt,right=4pt,top=2pt,bottom=2pt]
\begin{lstlisting}[basicstyle=\ttfamily\scriptsize,frame=none,backgroundcolor=\color{white}]
E-SEM-009 at 8:1: Function 'sign' may reach its end without returning a value of type int
    hint: Add 'return <int value>;' at the end of the function.
E-SEM-002 at 12:9: 'x' is already declared in this scope (line 11)
    hint: Use a different name, or assign to the existing variable with 'x = ...;'.
\end{lstlisting}
\end{tcolorbox}
\subsection*{r01\_division\_by\_zero.nova --- Runtime error: division by zero}
\noindent Compiles fine; the VM traps at run time. The optimizer must not fold this away.
\begin{lstlisting}[language=nova,]
fn main() {
    let a: int = 10;
    let zero: int = 5 - 5;
    print(a / 2);
    print(a / zero);
    print(1);
}
\end{lstlisting}
\begin{tcolorbox}[colback=green!4,colframe=green!45!black,title={Program output (optimized VM)},fonttitle=\bfseries\small,boxsep=2pt,left=4pt,right=4pt,top=2pt,bottom=2pt]
\begin{lstlisting}[basicstyle=\ttfamily\scriptsize,frame=none,backgroundcolor=\color{white}]
5
[runtime error: division by zero at line 8]
\end{lstlisting}
\end{tcolorbox}
\subsection*{r02\_stack\_overflow.nova --- Runtime error: stack overflow}
\noindent Unbounded recursion is stopped by the VM call-depth limit.
\begin{lstlisting}[language=nova,]
fn forever(n: int) -> int {
    return forever(n + 1) + 1;
}

fn main() {
    print(forever(0));
}
\end{lstlisting}
\begin{tcolorbox}[colback=green!4,colframe=green!45!black,title={Program output (optimized VM)},fonttitle=\bfseries\small,boxsep=2pt,left=4pt,right=4pt,top=2pt,bottom=2pt]
\begin{lstlisting}[basicstyle=\ttfamily\scriptsize,frame=none,backgroundcolor=\color{white}]

[runtime error: stack overflow: call depth exceeded 512 at line 5]
\end{lstlisting}
\end{tcolorbox}
\subsection*{r03\_index\_out\_of\_bounds.nova --- Runtime error: index out of bounds}
\noindent Array accesses are bounds-checked in the VM and in native code.
\begin{lstlisting}[language=nova,]
fn main() {
    let a: int[] = [1, 2, 3];
    print(a[2]);
    print(a[3]);
}
\end{lstlisting}
\begin{tcolorbox}[colback=green!4,colframe=green!45!black,title={Program output (optimized VM)},fonttitle=\bfseries\small,boxsep=2pt,left=4pt,right=4pt,top=2pt,bottom=2pt]
\begin{lstlisting}[basicstyle=\ttfamily\scriptsize,frame=none,backgroundcolor=\color{white}]
3
[runtime error: index out of bounds at line 7]
\end{lstlisting}
\end{tcolorbox}


\chapter{Benchmark raw output}\label{app:bench}
Raw output of \code{scripts/bench.sh} from the run behind Chapter~\ref{ch:perf}:

\begin{lstlisting}[]
benchmark  noregalloc_ms  regalloc_ms   cpython_ms   vs_noreg vs_cpython
fib35                74           54          654      1.37x      12.1x
primes               32           33          388       .96x      11.7x

benchmark     noinline_ms    inline_ms   cpython_ms    speedup
inline                103           47          n/a      2.19x
\end{lstlisting}


\chapter{Specifications of the VM instruction set}
The complete instruction set is listed in Chapter~\ref{ch:backends} (generated from the compiler); the language grammar is Listing~\ref{lst:grammar}; the SQL statement summary is Listing~\ref{lst:sqlgrammar}.
